% latexmk -pdf -pvc main.tex
\documentclass{../util/zariski-small}


\title{Synthetic Adequate Modules}
\author{Felix Cherubini, J. Dan Christensen, Fabian Endres, Thomas Thorbjørnsen}
\date{\today}

% --- Macros --- %
\newcommand{\earlyqed}{\hfill$\square$}
\newcommand{\lra}       {\longrightarrow}
\newcommand{\llra}[1]   {\stackrel{#1}{\lra}}  % labelled long right arrow

\newcommand{\hood}[2]{\mathcal{N}_{#2}(#1)}

\renewcommand{\fib}{\mathsf{fib}}
\newcommand{\susp}{\mathsf{\Sigma}}
\newcommand{\ap}{\mathsf{ap}}
\newcommand{\UU}{\mathcal{U}}
\renewcommand{\Prop}{\mathsf{Prop}}
\renewcommand{\refl}{\mathsf{refl}}
\renewcommand{\id}{\mathsf{id}}
\newcommand{\pr}{\mathsf{pr}}
\newcommand{\tot}{\mathsf{tot}}
\newcommand{\defeq}{:\equiv}
\newcommand{\sbt}{{\,\begin{picture}(-1,1)(-1,-3)\circle*{3}\end{picture}\ } }
\newcommand{\pto}{\to_{\sbt}}
\newcommand{\IsEquiv}{\mathsf{IsEquiv}}

\newcommand{\mcal}[1]{\mathcal{#1}} % Shorthand for calligraphic font
\newcommand{\mcalbf}[1]{\mathbf{\mcal{#1}}} % Caligraphic bold math
\newcommand{\op}{\textup{op}} % Opposite category
\newcommand{\sigmatype}[1]{\underset{#1}{\Sigma}} % Sigma type
\newcommand{\pitype}[1]{\underset{#1}{\Pi}} % Pi type
\newcommand{\argmhole}{\ \cdot\ } % Hole
\newcommand{\set}[1]{\{#1\}} % Sets
\newcommand{\setbuild}[2]{\{#1 \mid #2\}} % Set-builder
\newcommand{\ideal}[1]{\langle #1 \rangle}
\newcommand{\quot}[2]{#1/#2}
\newcommand{\tangent}[2]{\mathrm{T}_{#2}#1}
\newcommand{\der}[2]{d#1_#2}
\newcommand{\can}{\mathrm{can}}
\newcommand{\inl}{\mathrm{inl}}
\newcommand{\inr}{\mathrm{inr}}
\newcommand{\inc}{\mathrm{inc}}
\newcommand{\ev}{\mathrm{ev}}

\newcommand{\sSet}{\textup{sSet}}
\newcommand{\Kan}{\textup{Kan}}
\newcommand{\Ab}{\textup{Ab}}
\newcommand{\modcat}[1]{\textup{Mod}_{#1}} % Category of R-modules
\newcommand{\finfree}{\modcat{R}^{\mathrm{fin.free}}}
\newcommand{\adeq}{\modcat{R}^{\mathrm{adeq}}} % Category of adequate $R$-modules
\newcommand{\finadeq}{\modcat{R}^{\mathrm{fin.adeq}}} % Category of finitely adequate
\newcommand{\LEM}{\textup{LEM}}
\newcommand{\AC}{\textup{AC}}
\newcommand{\pType}{\Type_\ast}

% --- Redeclare math operators --- %
% https://tex.stackexchange.com/questions/175251/how-to-redefine-a-command-using-declaremathoperator
\makeatletter
\newcommand\RedeclareMathOperator{%
  \@ifstar{\def\rmo@s{m}\rmo@redeclare}{\def\rmo@s{o}\rmo@redeclare}%
}
% this is taken from \renew@command
\newcommand\rmo@redeclare[2]{%
  \begingroup \escapechar\m@ne\xdef\@gtempa{{\string#1}}\endgroup
  \expandafter\@ifundefined\@gtempa
     {\@latex@error{\noexpand#1undefined}\@ehc}%
     \relax
  \expandafter\rmo@declmathop\rmo@s{#1}{#2}}
% This is just \@declmathop without \@ifdefinable
\newcommand\rmo@declmathop[3]{%
  \DeclareRobustCommand{#2}{\qopname\newmcodes@#1{#3}}%
}
\@onlypreamble\RedeclareMathOperator
\makeatother

% \DeclareMathOperator{\Hom}{Hom}
% \DeclareMathOperator{\Ext}{Ext}
\DeclareMathOperator{\Nat}{Nat}
\DeclareMathOperator{\Ker}{Ker}
\DeclareMathOperator{\Cok}{Cok}
\RedeclareMathOperator{\Im}{Im}
\DeclareMathOperator{\Coim}{Coim}


\begin{document}

\maketitle

\begin{abstract}

\end{abstract}

\tableofcontents

\section*{Introduction}
\input{overview.tex}
\paragraph{Contribution and organization of the article.}
The first motivating question behind this article is whether synthetic algebraic geometry admits a notion similar to the coherent sheaves in classical algebraic geometry. Coherent sheaves share properties with finite dimensional vector spaces and appear for example in Serre duality, an important tool of classical algebraic geometry.

Before going into the challenges with finding a synthetic analogue of coherent sheaves, we will give some background.
Let $R$ be the base ring of synthetic algebraic geometry, as mentioned above. 
The three axioms of synthetic algebraic geometry imply that some properties of $R$ are close to that of a field.
For example, it was noticed in~\cite{diffgeo-article} that any finitely presented $R$-module $M$ is reflexive, i.e.\ canonically isomorphic to its double dual $M^{\ast\ast}$ where $M^\ast$ is $\Hom_R(M,R)$. In this article, we define \emph{finitely adequate $R$-modules} to be $R$-modules, which are merely the kernel of a linear map between finitely presented $R$-modules. Reflexivity extends to this category of $R$-modules.

In synthetic algebraic geometry, we work with a generalization of $\mathcal O_X$-module sheaves called $R$-module \emph{bundles}.
Module bundles are just maps $X\to \Mod{R}$, assigning to each point of a type $X$ an $R$-module. In \cite{draft}, the abelian category of \emph{weakly quasi-coherent} $R$-modules was defined and bundles of such modules on schemes were used for cohomological calculations. While these bundles are more general than the usual quasi-coherent $\mathcal O_X$-modules, it was still possible to adapt many proofs from the literature. Similarly, a synthetic replacement for coherent sheaves should admit the same arguments used in known proofs. Any direct adaption of the definition of coherent sheaves is bound to fail: Classically, it only works well if finitely presented modules over the base ring form an abelian category, which is known to be false for $R$.

We will now give some requirements for a basic proof of Serre duality, like it is found in \cite{Hartshorne}. Let $X$ be a smooth projective scheme. We would ideally need one subcategory $\mathcal A$ of the $R$-module bundles on $X$ and one subcategory $\mathcal A^\prime\subseteq \Mod{R}$ with the following properties:

\begin{enumerate}[(i)]
\item Locally free sheaves are contained in $\mathcal A$.
\item For any $\mathcal F$ in $\mathcal A$ all cohomology modules $H^i(X,\mathcal F)$ are in $\mathcal A^\prime$.
\item The dualization functor given by $\Hom_R(\argmhole,R)$ is exact on $\mathcal{A}^\prime$.
\end{enumerate}

Our initial hope was that $\mathcal A^\prime$ can be the category of \emph{finitely adequate $R$-modules} and $\mathcal A$ can be chosen as bundles on $X$ with values in $\mathcal A^\prime$. In work in progress based on this article it turned out that $\mathcal A$ has to be a smaller category, for example the category of vector bundles on $X$, in order for the cohomology to take values in finitely adequate $R$-modules. With this restriction it seems to be the case that Serre duality follows with a proof similar to that of Hartshorne. The externalization of this proof would give a constructive proof of Serre duality over a general ring. We do not know of an existing constructive proof or any proof in this generality.

Beyond what is needed for Serre duality there are many good and surprising properties of the abelian category of finitely adequate $R$-modules:

\begin{itemize}
\item The module $\Hom_R(M, N)$ is finitely adequate when $M$ and $N$ are.
\item Dualization is an equivalence $\Hom_R(\argmhole, R) : (\finadeq)^{\op} \to \finadeq$.
\item The injective modules are exactly the finitely presented modules and the projective modules are exactly the finitely copresented modules.
\item Finitely adequate $R$-modules are closed under extensions (inside $\Mod{R}$).
\item The category $\finadeq$ of finitely adequate modules is closed symmetric monoidal, with tensor product the double dual of the usual tensor product.
\end{itemize}

All these results are shown in \cref{adequate-section}, together with the theorem that finitely adequate $R$-modules form an abelian category. Already the prior work in \cite{diffgeo-article} used a differential geometry argument to show projectivity of finitely copresented $R$-modules. The differential geometry of the underlying type of  some modules plays even more a role in the present work and is essential in our proof that finitely adequate $R$-modules are closed under extension.

In \cref{infinitesimally-linear-section}, as a preparation for the central proofs in \cref{adequate-section}, we study how the differential geometric properties of the underlying type relates to the algebraic structure of a $R$-module.
The synthetic definition of tangent spaces works for general types. For \emph{infinitesimally linear} types, these tangent spaces are $R$-modules. If $M$ is an infinitesimally linear $R$-module, we can compare $M$ with its own tangent space via a canonical map, which always exists and which we prove to be linear. If, for a general $R$-module, this canonical map turns out to be an equivalence, we call the module \emph{tangential}. We show that tangential $R$-modules are always infinitesimally linear and that the canonical map is an isomorphism to their tangent space at 0.  To our surprise, many $R$-modules in synthetic algebraic geometry are tangential, in particular weakly quasi-coherent and finitely adequate $R$-modules.

An alternative to the explicit definition of finitely adequate $R$-modules given above would be to describe the finitely adequate $R$-modules as the smallest abelian subcategory of $\Mod{R}$ containing all finite free $R$-modules. More surprisingly, finitely adequate $R$-modules are also described by the construction of the free abelian category in \cite{adelman-construction}. For this reason, we considered calling finitely adequate $R$-modules ``Adelman modules'', but decided to name them after the existing notion of adequate modules \cite[Chapter 46]{stacks-project}, once we were aware of it.
\rednote{``quite similar'':  should we say that it \emph{is} the free abelian category on the finite free modules?}

Apart from the free abelian category and slight variations on the definition given in the beginning, we also asked if one of the following categories coincides with the finitely adequate $R$-modules:
\begin{center}
\begin{enumerate}[(i)]
\item Reflexive modules, i.e.\ all $R$-modules $M$ with $M^{\ast\ast}=M$ by the canonical isomorphism.
\item $R$-modules $M$ that are not not finite free.
\end{enumerate}
\end{center}
The answer turned out to be negative in both cases. We present these results in \cref{examples-section} together with other examples that show the relations of the various types of $R$-modules which appear in synthetic algebraic geometry.

\paragraph{Acknowledgements.}
\rednote{TODO: Add everything that should be added.}

The work in \Cref{adequate-section} profited from discussions with Mohamed Barakat and Marc Nieper-Wißkirchen.

We thank David Wärn, Ingo Blechschmidt and Marc Nieper-Wißkirchen for a fruitful discussion at the Dagstuhl-Seminar 26121. At the meeting, David Wärn explained his idea for proving that finitely adequate modules are stable under extension, which we use in this article.
Work on this article by Felix Cherubini was supported by the ForCUTT project, ERC advanced grant 101053291.



\section{Infinitesimally Linear Types}
In this section, we build on the theory of tangent spaces developed by Cherubini et al.~\cite{diffgeo-article}.
We introduce and compare two coexisting $R$-module structures on the tangent space of an infinitesimally linear type: the pointwise structure and the microlinear structure.
Finally, we conclude the section with a linearization principle: we show that under suitable conditions on the modules, any pointed lifting problem in a commutative square can be solved linearly, provided a non-linear pointed solution exists.

\subsection{Infinitesimal Types and Tangent Spaces}

In synthetic algebraic geometry, tangent spaces are defined in terms of pointed maps out of the first-order neighbourhood of $0$ in the ring $R$.
This subsection introduces the concept of $n$-th order neighbourhoods around a point, which we use to recall the definition of tangent spaces from \cite{diffgeo-article}.

\begin{definition}\label{def:hood-res}
  Let $(X,x)$ be a pointed type.
  The \notion{$n$-th order neighbourhood around $x$} is the subtype $\hood{X,x}{n} \defeq \setbuild{x' : X}{\exists I \substack{\subseteq \\ \mathrm{fg.ideal}} R, I^{n+1} = 0, I = 0 \rightarrow x' = x}$.
\end{definition}

Taking $n$-th order neighbourhoods defines a functor $\hood{\argmhole}{n} : \pType \rightarrow \pType$.
The action on maps $f : (X,x) \rightarrow_\ast (Y,y)$ is given by $\hood{f}{n} \defeq f\mid_{\hood{X,x}{n}}$.

\begin{proposition}\label{prp:hood-fun-welldef}
  $\hood{\argmhole}{n}$ is a well-defined functor.
\end{proposition}

\begin{proof}
  We need to show that for any $x' \in \hood{X,x}{n}$, we have $f(x') \in \hood{Y,y}{n}$.
  By assumption, we merely have a finitely generated ideal $I \subseteq R$ such that $I^{n+1} = 0$ and if $I = 0$ then $x' = x$.
  It follows that $f(x') \in \hood{Y,y}{n}$, as if $I = 0$, then $f(x') = f(x) = y$ by applying $f$ to $x' = x$ and by $f$ being pointed.
\end{proof}

Observe that we have an inclusion $i_x : \hood{X,x}{n} \hookrightarrow X$ natural in $(X,x)$.

\begin{definition}\label{def:hood}
  A type $X$ is an \notion{$n$-th order neighbourhood around $x : X$} if the inclusion $i_x : \hood{X,x}{n} \hookrightarrow X$ is an equivalence of types.
\end{definition}

\begin{example}\label{ex:R[x]/(x^{n+1})}
  Consider the type $X \defeq \Spec(\quot{R[x_1, \dots, x_m]}{\ideal{x_1, \dots, x_m}^{n+1}})$.
  The ideal $\ideal{x_1, \dots, x_m}^{n+1}$ is generated by all monomials of total degree $n+1$.
  A point $f : X$ is an $R$-algebra homomorphism $f : \quot{R[x_1, \dots, x_m]}{\ideal{x_1, \dots, x_m}^{n+1}} \to R$, which is determined by the values $f(x_i) = r_i \in R$ for $i \in \set{1,\dots,m}$.
  The imposed relations require that for any monomial $p(x_1, \dots, x_m)$ of degree $n+1$, $f(p(x_1, \dots, x_m)) = p(r_1, \dots, r_m) = 0$.
  Equivalently, the ideal $I \defeq \ideal{r_1, \dots, r_m} \subseteq R$ satisfies $I^{n+1} = 0$.

  The type $X$ is an $n$-th order neighbourhood around $0$, i.e. the map $i : \hood{X,0}{n} \hookrightarrow X$ is an equivalence.
  To see this, consider $f : X$ given by $(r_1, \dots, r_m)$.
  Let $I \defeq \ideal{r_1, \dots, r_m}$. Then $I^{n+1} = 0$ by construction.
  If $I = 0$, then each $r_i=0$, implying that $f = 0$.
\end{example}

For the case $n=1$ in the preceding example, we introduce a special notation.
We define the \notion{$m$-dimensional first-order infinitesimal disk}, or just the \notion{$m$-disk}, as the type $\D(m) \defeq \Spec(\quot{R[x_1,\dots, x_m]}{\ideal{x_1,\dots,x_m}^2})$.
By duality, we can identify $\D(m)$ with the type $\{ x : R^m \mid x_ix_j = 0 \text{ for all } i,j \}$.
The $m$-disk is a first-order neighbourhood around $0$.
Thus, whenever we consider $\D(m)$ as a pointed type, we will always mean the pointed type $(\D(m), 0)$.

Let $(X,x)$ and $(Y,y)$ be pointed types.
We introduce the notation $(Y,y)^{(X,x)}$ for the type of pointed maps $(X,x) \rightarrow_\ast (Y,y)$.
If the pointing on $X$ is understood from context, then we will simply write $(Y,y)^X$.

\begin{definition}\label{def:tangent}
  Let $(X,x)$ be a pointed type.
  The tangent space of $X$ at $x$ is the type $\tangent{X}{x} \defeq (X,x)^{\D(1)}$.
\end{definition}

The tangent space construction defines a functor $\tangent{}{} : \pType \rightarrow \pType$.
If $f : (X,x) \rightarrow_\ast (Y,y)$ is a pointed map, the derivative of $f$ is the map $\der{f}{x} \defeq f_\ast : \tangent{X}{x} \rightarrow_\ast \tangent{Y}{y}$ induced by post-composition.
Taking tangent spaces is a local construction in the following sense.

\begin{lemma}\label{prp:hood-factorization}
  Let $(X,x)$ and $(Y,y)$ be pointed types.
  If $X$ is an $n$-th order neighbourhood around $x$, then every pointed map $f : (X,x) \rightarrow_\ast (Y,y)$ factors uniquely through $\hood{Y,y}{n}$.
\end{lemma}

\begin{proof}
  This is immediate by functoriality of $\hood{\argmhole}{n}$, \Cref{prp:hood-fun-welldef}, as $f$ factors into $i_y \circ \hood{f}{n} \circ i_x^{-1}$.
  This factorization is unique, since $i_y$ is an inclusion.
\end{proof}

\begin{remark}\label{rem:tangent-factorization}
  Let $(X,x)$ be a pointed type.
  Specializing \Cref{prp:hood-factorization} to $\tangent{X}{x} \defeq (X,x)^{\D(1)}$, we observe that any tangent vector $\nu : \tangent{X}{x}$ has its image entirely in $\hood{X,x}{1}$.
  Thus, $\tangent{X}{x} = \tangent{\hood{X,x}{1}}{x}$.
\end{remark}

% Note that $\D(1) = \Spec(\quot{R[x]}{\ideal{x}^2})$ can be identified with $\sum_{x:R} x^2 = 0$, namely the square-zero elements of $R$.
% We make use of this in the following construction.

\subsection{Module Structures on Tangent Spaces}

Let $(X,x)$ be a pointed type.
If $X$ is an $R$-module and $x = 0$, then we will describe the pointwise $R$-module structure on the tangent space $\tangent{X}{x}$.
In the case that $(X,x)$ is infinitesimally linear (to be defined in \Cref{def:inf-lin}), we can equip $\tangent{X}{x}$ with a second $R$-module structure.
In this subsection, we will compare these module structures when $M$ is an $R$-module such that $(M,0)$ is infinitesimally linear.


We begin by defining the pointwise module structure on the tangent space $\tangent{M}{0}$ of an $R$-module $M$.
Let $\upsilon, \nu$ be tangent vectors in $\tangent{M}{0}$.
The pointwise addition is given by $(\upsilon + \nu)(\varepsilon) = \upsilon (\varepsilon) + \nu (\varepsilon)$ for all $\varepsilon : \D(1)$.
The scalar multiplication is given by $(r \cdot \upsilon)(\varepsilon) = r \cdot \upsilon (\varepsilon)$ for all $r \in R$ and $\varepsilon : \D(1)$.
It is clear that these operations are well-defined and define a module structure on $\tangent{M}{0}$.

\begin{remark}
  Let $M$ and $N$ be $R$-modules.
  Taking the derivative of a pointed map $f : M \to_\ast N$ yields a pointed map $\der{f}{0} : \tangent{M}{0} \to_\ast \tangent{N}{0}$.
  It is only when $f$ is linear that $\der{f}{0}$ is linear with respect to the pointwise module structures.
\end{remark}

We define the \notion{canonical map} $\can : M \rightarrow \tangent{M}{0}$ by $m \mapsto (\varepsilon \mapsto \varepsilon \cdot m)$.
It is clear that $\can$ is $R$-linear with respect to the pointwise module structure and natural with respect to $R$-linear maps.

So far, we have only considered the tangent spaces of types that are already $R$-modules.
However, the tangent space $\tangent{X}{x}$ of a general pointed type $(X,x)$ can also carry a natural $R$-module structure.
The condition under which this occurs, known as \notion{infinitesimal linearity}, is formulated using certain canonical maps between infinitesimal disks.

For any $m,n : \N$, we define the following inclusion:
\begin{align*}
  \inc_{m,n} : \D(m) \vee \D(n) & \rightarrow \D(m+n) \\
  \inl (\varepsilon) & \mapsto (\varepsilon, 0) \\
  \inr (\varepsilon) & \mapsto (0, \varepsilon)
\end{align*}
The diagonal of the $1$-disk is defined as:
\begin{align*}
  D : \D(1) & \rightarrow \D(2) \\
  \varepsilon & \mapsto (\varepsilon,\varepsilon)
\end{align*}

\begin{definition}\label{def:inf-lin}
  Let $(X,x)$ be a pointed type.
  $(X,x)$ is \notion{infinitesimally linear} if for every $m, n : \N$ the induced map
  \begin{align*}
    \inc_{m,n}^* : (X,x)^{\D (m+n)} \rightarrow (X,x)^{\D(m) \vee \D(n)}
  \end{align*}
  is an equivalence.
\end{definition}

Infinitesimal linearity is a local property in the following sense:

\begin{proposition}\label{prp:inf-lin-hood}
  Let $(X,x)$ be a pointed type.
  $X$ is infinitesimally linear at $x$ if and only if $\hood{X,x}{1}$ is infinitesimally linear at $x$.
\end{proposition}

\begin{proof}
  By \Cref{prp:hood-factorization} any map $\D(n) \rightarrow_\ast (X,x)$ factors uniquely through $\hood{X,x}{1}$.
  Thus, we have the desired lifting property for $(X,x)$ if and only if we have it for $\hood{X,x}{1}$.
\end{proof}

Let $(X,x)$ be infinitesimally linear.
The tangent space $\tangent{X}{x}$ carries an $R$-module structure, which we call the microlinear module structure.
Let $\upsilon$ and $\nu$ be tangent vectors in $\tangent{X}{x}$.
The microlinear addition $\upsilon \boxplus \nu$ is defined by the composite $(\inc_{1,1}^\ast)^{-1}[\upsilon,\nu] \circ D$.
In other words, we obtain this map by lifting in the following diagram
\begin{center}
  \begin{tikzcd}
    & \D(1) \vee \D(1) \ar{rd}{[\upsilon,\nu]} \ar{d}{\inc_{1,1}} \\
    \D(1) \ar{r}{D} \ar[bend right]{rr}{\upsilon \boxplus \nu} & \D(2) \ar[dashed]{r}{\exists !} & X
  \end{tikzcd}
\end{center}
The microlinear scalar multiplication is defined by $(r \boxdot \upsilon)(\varepsilon) = \upsilon(r\varepsilon)$ for all $r : R$ and $\varepsilon : \D(1)$.
Indeed, to prove that this defines a module structure, one can adapt the proof of \cite[Theorem 4.2.19]{david-orbifolds}.
We will give a direct proof.

\begin{theorem}\label{thm:tangent-inf-lin}
  Let $(X,x)$ be infinitesimally linear.
  Then the microlinear module structure on $\tangent{X}{x}$ is well-defined. \earlyqed
\end{theorem}

\begin{proof}
  We need to verify the axioms to have an $R$-module structure.
  
  Unitality is clear by definition.
  
  Associativity follows from infinitesimal linearity.
  Let $\phi : \D(m) \to X$ and $\psi : \D(n) \to X$ be two different sums of tangent vectors.
  We let $\phi_i$ and $\psi_j$ denote the restrictions to the $i$-th and $j$-th components of $\phi$ and $\psi$, respectively.
  The sum $\phi \boxplus \psi : \D(m+n) \to X$ is uniquely given by $[\phi_1,\dots,\phi_m,\psi_1,\dots,\psi_n]$.
  Thus $\phi \boxplus \psi = \phi_1 \boxplus \dots \boxplus \phi_m \boxplus \psi_1 \boxplus \dots \boxplus \psi_n$, as desired.

  Distributivity is also clear.
  Let $\upsilon$ and $\nu$ be tangent vectors in $\tangent{X}{x}$, and $r$ be an arbitrary element of $R$.
  Then $((r\boxdot\upsilon)\boxplus(r\boxdot\nu))(\varepsilon)$ is given by $[r \boxdot \upsilon, r \boxdot \nu]$.
  Likewise, $(r \boxdot (\upsilon \boxplus \nu))(\varepsilon)$ evaluates to $(\upsilon \boxplus \nu)(r\varepsilon)$ for each $\varepsilon$ in $\D(1)$.
  Thus the second tangent vector is also given by $[r \boxdot \upsilon, r \boxdot \nu]$, as desired.
\end{proof}

The microlinear module structure is natural in the sense that the derivative of any pointed map between infinitesimally linear types is linear.
The following proposition generalizes \cite[Lemma 2.2.6]{diffgeo-article} to this setting.

\begin{proposition}\label{prp:derivative-inf-lin}
  Let $f : (X,x) \rightarrow_\ast (Y,y)$ be a pointed map between infinitesimally linear types.
  Then $\der{f}{x}$ is $R$-linear with respect to the microlinear module structure.
\end{proposition}

\begin{proof}
  Let $\upsilon$ be a tangent vector in $\tangent{X}{x}$.
  Then we have $f(r\boxdot\upsilon)(\varepsilon) = f(\upsilon)(r\varepsilon) = (r \boxdot f(\upsilon))(\varepsilon)$ for each $\varepsilon$ in $\D(1)$.
  Thus, $\der{f}{x}(r\boxdot\upsilon) = r \boxdot \der{f}{x}(\upsilon)$, as desired.

  Let $\upsilon$ and $\nu$ be tangent vectors in $\tangent{X}{x}$.
  Then both triangles in the diagram
  \begin{center}
    \begin{tikzcd}
      & \D(1) \vee \D(1) \ar{rd}{[f(\upsilon),f(\nu)]} \ar{d}{\inc_{1,1}} \\
      \D(1) \ar{r}{D} \ar[bend right]{rr}{f(\upsilon \boxplus \nu)} & \D(2) \ar[dashed]{r}{\exists !} & X
    \end{tikzcd}
  \end{center}
  commute, because $f \circ [\upsilon,\nu] = [f(\upsilon),f(\nu)]$.
  Thus $\der{f}{x}(\upsilon \boxplus \nu) = \der{f}{x}(\upsilon) \boxplus \der{f}{x}(\nu)$, as desired.
\end{proof}

If $M$ is an infinitesimally linear $R$-module, its tangent space $\tangent{M}{0}$ carries two potentially different $R$-module structures: the pointwise structure and the microlinear structure. We now compare these two structures.

\begin{theorem}\label{thm:tangent-ptw-micro-add}
  Suppose that $M$ is an infinitesimally linear $R$-module.
  Then the pointwise and microlinear addition on $\tangent{M}{0}$ are the same.
\end{theorem}

\begin{proof}
  Let $\upsilon,\nu : \tangent{M}{0}$.
  The pointwise addition $\upsilon + \nu$ factors through $\D(2)$ by $(\varepsilon,\delta) \mapsto \upsilon(\varepsilon) + \nu(\delta)$.
  It is clear that this map is given by $[\upsilon,\nu]$.
  Thus $\upsilon + \nu = \upsilon \boxplus \nu$.
\end{proof}

\begin{remark}\label{rem:ptw-micro-scalar}
  In general, the scalar multiplication given by the pointwise and microlinear module structures differs.
  Let $\nu : \tangent{M}{0}$ be a tangent vector on $M$, $r : R$, and $\varepsilon : \D(1)$.
  Then $(r\boxdot \nu)(\varepsilon) = \nu(r\varepsilon)$ and $(r \cdot \nu)(\varepsilon) = r\nu(\varepsilon)$.
  We can observe that $r\boxdot\nu = r\cdot \nu$ holds if and only if $\nu$ respects scalar multiplication.
\end{remark}

This comparison allows us to show that the canonical map is also linear with respect to the microlinear structure:

\begin{corollary}\label{cor:can-nat-inf-lin}
  Let $(M,0)$ be an infinitesimally linear $R$-module.
  Then the map ${\can : M \rightarrow \tangent{M}{0}}$ is $R$-linear with respect to the microlinear module structure.
\end{corollary}

\begin{proof}
  The map $\can$ is $R$-linear with respect to the pointwise module structure on $\tangent{M}{0}$.

  Firstly, $\can$ respects the microlinear addition.
  To see this, let $\upsilon$ and $\nu$ be tangent vectors in $\tangent{M}{0}$.
  By \Cref{thm:tangent-ptw-micro-add}, we have $\can(\upsilon \boxplus \nu) = \can(\upsilon + \nu) = \can(\upsilon) + \can(\nu) + \can(\upsilon) \boxplus \can(\nu)$. 

  Secondly, $\can$ respects the microlinear scalar multiplication.
  Let $m : M$, $\varepsilon : \mathbb{D}(1)$ and $r : R$.
  Then we have $(r\boxdot\mathrm{can}(m))(\varepsilon) = \mathrm{can}(m)(r\varepsilon) = r\varepsilon m = \varepsilon rm = \mathrm{can}(rm)(\varepsilon)$.
\end{proof}

\begin{remark}
  The pointwise and microlinear module structures on $\tangent{M}{0}$ agree if $\can$ is surjective.
\end{remark}

\subsection{Tangential Modules}

In this subsection, we study tangential modules (\Cref{def:tangential}).
We show that this class of modules is closed under standard operations like finite limits and finite colimits.
We then conclude this subsection by showing that pointed maps between such modules admit a linearization.

\begin{definition}\label{def:tangential}
  An $R$-module $M$ is \notion{tangential} if $\can : M \rightarrow \tangent{M}{0}$ is an equivalence.
\end{definition}

We can generalize $\can$ to a family of canonical maps $\can^k : M^k \to (M,0)^{\D(k)}$ for all $k \geq 1$, defined by
\[
\can^k (m_1, \dots, m_k) (\varepsilon_1, \dots, \varepsilon_k) = \sum_{i=1}^k \varepsilon_i m_i .
\]

\begin{proposition}\label{prp:tangential-implies-gen-can-equiv}
  Let $M$ be an $R$-module. If $M$ is tangential, then $\can^k$ is an equivalence for all $k \geq 1$.
\end{proposition}

\begin{proof}
  We will show this by induction.
  The base case $k=1$ is trivial by assumption.
  Now, assume the statement holds for $k \ge 1$.
  Let $f : \D(k+1) \to M$ be a pointed map.
  Since $M$ is tangential, the map $\delta \mapsto f(0, \dots, 0, \delta)$ is of the form $\delta m_{k+1}$ for a unique $m_{k+1} : M$.
  Consider the pointed map $\hat{f}(\varepsilon_1, \dots, \varepsilon_k, \delta) \defeq f(\varepsilon_1, \dots, \varepsilon_k, \delta) - \delta m_{k+1}$, which satisfies $\hat{f}(0, \dots, 0, \delta) = 0$.
  For any $\delta : \D(1)$, the map $(\varepsilon_1, \dots, \varepsilon_k) \mapsto \hat{f}(\varepsilon_1, \dots, \varepsilon_k, \delta)$ is a pointed map $\D(k) \to M$.
  By the induction hypothesis, this map is uniquely of the form $\sum_{i=1}^k \varepsilon_i g_i(\delta)$ for some functions $g_i : \D(1) \to M$.
  Define $\hat{g_i}(\delta) \defeq g_i(\delta) - m_i$ where $m_i \defeq g_i(0)$.
  Since $M$ is tangential, the maps $\hat{g_i} : \D(1) \to M$ are uniquely of the form $\delta m'_i$ for some $m'_i : M$.
  Thus we have $g_i(\delta) = m_i + \hat{g_i}(\delta) = m_i + \delta m_i'$.
  We can then compute:
  \[ \hat{f}(\varepsilon_1, \dots, \varepsilon_k, \delta) = \sum_{i=1}^k \varepsilon_i (m_i + \delta m'_i) = \sum_{i=1}^k \varepsilon_i m_i + \sum_{i=1}^k \varepsilon_i \delta m'_i. \]
  Since $(\varepsilon_1, \dots, \varepsilon_k, \delta) : \D(k+1)$, we have that $\varepsilon_i \delta = 0$ for all $i$.
  Therefore, the second term vanishes, yielding $\hat{f}(\varepsilon_1, \dots, \varepsilon_k, \delta) = \sum_{i=1}^k \varepsilon_i m_i$, which shows that $f(\varepsilon_1, \dots, \varepsilon_k, \delta) = \sum_{i=1}^{k+1} \varepsilon_i m_i$ for unique $m_1, \dots, m_{k+1} : M$.
  Thus $\can^{k+1}$ is an equivalence.
\end{proof}

\begin{corollary}\label{cor:tangential-inf-lin}
  Let $M$ be a tangential $R$-module. Then $(M,0)$ is infinitesimally linear.
\end{corollary}

\begin{proof}
  We have the following chain of equivalences factoring $\inc_{m,n}$:
  \begin{center}
    \begin{tikzcd}
      M^{\D(m+n)} \ar{r}{\inc_{m,n}} & M^{\D(m)\vee\D(n)} \ar{rd}[right, below]{\simeq} \\
      & & M^{\D(m)}\times M^{\D(n)} \\
      M^{m+n} \ar{uu}{\can^{m+n}}[right]{\simeq} \ar{r}[below]{\simeq} & M^m \times M^n \ar{ru}{\can^m \times \can^n}[right, below]{\simeq}
    \end{tikzcd}
  \end{center}
  The maps $\can^{m+n}$ and $\can^m \times \can^n$ are equivalences by \Cref{prp:tangential-implies-gen-can-equiv}.
\end{proof}

The properties of being tangential and infinitesimally linear are closed under many common categorical constructions.
This will be important to prove that various modules are themselves tangential.
For instance, we will use this in \Cref{ex:power-series-tangential} to show that the polynomial ring $R[x]$ and power series ring $R[[x]]$ are tangential.

\begin{proposition}[{{\cite[Proposition 5.0.1]{david-orbifolds}}}]\label{prp:lim-inf-lin}
  Infinitesimally linear types are closed under limits.
  If $(X,x)$ is infinitesimally linear, then $(X^A, \mathrm{const}_x)$ is infinitesimally linear for every type $A$.
\end{proposition}

\begin{proposition}\label{prp:exp-tangential}
  Let $M$ be a tangential $R$-module.
  Then for any type $U$, the $R$-module $M^U$ is tangential.
\end{proposition}

\begin{proof}
  % One of the mapping spaces is pointed, so there's a tiny thing to check.
  It is straightforward to check that $\tangent{(M^U)}{0} = (\tangent{M}{0})^U$.
  The latter type is equivalent to $M^U$ by assumption,
  and since the module structure on $M^U$ is pointwise, it follows that the
  resulting equivalence $M^U \to \tangent{(M^U)}{0}$ is homotopic to the canonical map.
\end{proof}

To show that tangential and infinitesimally linear $R$-modules are closed under finite limits and finite colimits, we need some lemmas regarding projectivity.

\begin{lemma}\label{lem:disk-projective}
  For any $m : \N$ the types $\D(m)$ are projective among types.
\end{lemma}

\begin{proof}
  Let $f : X \to \D(m)$ be a surjection. The family $\fib_f : \D(m) \to \UU$ is inhabited by assumption. By Zariski local choice, there exists a unimodular vector $f_1, \dots, f_n : \quot{R[x_1,\dots,x_m]}{\langle x_1, \dots, x_m \rangle^2}$ such that each $f_i$ admits a local section $s_i : D(f_i) \to X$. Since $\quot{R[x_1,\dots,x_m]}{\langle x_1, \dots, x_m \rangle^2}$ is a local ring, at least one of the $f_i$ is a unit, say $f_k$. Then $s_k$ is a global section of $f$.
\end{proof}

\begin{lemma}\label{lem:projective-exact}
  Let $(X,x)$ be projective among types.
  The functor $(\argmhole,0)^{(X,x)} : \modcat{R} \rightarrow \modcat{R}$ is exact.
\end{lemma}

\begin{proof}
  Let $M \twoheadrightarrow N$ be an epimorphism of $R$-modules, and suppose we have a pointed map $f : (X,x) \rightarrow_\ast (N,0)$.
  We merely get a lift $\hat{g} : X \rightarrow M$ of $f$ along the epi.
  The map $g = \hat{g} - \hat{g}(0)$ is a pointed lift of $f$, as desired.
\end{proof}

\begin{proposition}\label{prp:tangent-exact}
  The functors $(\argmhole,0)^{\D(m+n)},\ (\argmhole,0)^{\D(m) \vee \D(n)}$ are exact on $\modcat{R}$.
  In particular, $\tangent{}{0}$ is exact on $\modcat{R}$.
\end{proposition}

\begin{proof}
  Immediate by \Cref{lem:disk-projective} and \Cref{lem:projective-exact}.
\end{proof}

\begin{proposition}\label{prp:op-closure-inf-lin-tangential}
  Infinitesimally linear $R$-modules at $0$ are closed under finite colimits and extensions.
  Tangential $R$-modules are closed under finite limits, finite colimits and extensions.
\end{proposition}

\begin{proof}
  This follows from \Cref{prp:tangent-exact}.
  All cases follow the same pattern.
  We will show how to prove this for the case of infinitesimally linear $R$-modules. An analogous proof works for the other cases.
  The key idea is to use the $5$-lemma to deduce that $\inc_{m,n}^\ast$ is an equivalence, as illustrated in the diagrams below.
  \begin{center}
    \begin{tikzcd}
      A^{\D(m+n)} \ar[hook]{d} \ar{r}[above]{\simeq}[below]{\inc_{m,n}^\ast} & A^{\D(m)\vee\D(n)} \ar[hook]{d} \\
      B^{\D(m+n)} \ar{r}[above]{\simeq}[below]{\inc_{m,n}^\ast} \ar[two heads]{d} & B^{\D(m)\vee\D(n)} \ar[two heads]{d} \\
      C^{\D(m+n)} \ar{r}[above]{}[below]{\inc_{m,n}^\ast} & C^{\D(m)\vee\D(n)}
    \end{tikzcd}
    \qquad
    \begin{tikzcd}
      A^{\D(m+n)} \ar[hook]{d} \ar{r}[above]{\simeq}[below]{\inc_{m,n}^\ast} & A^{\D(m)\vee\D(n)} \ar[hook]{d} \\
      B^{\D(m+n)} \ar{r}[above]{}[below]{\inc_{m,n}^\ast} \ar[two heads]{d} & B^{\D(m)\vee\D(n)} \ar[two heads]{d} \\
      C^{\D(m+n)} \ar[]{r}[above]{\simeq}[below]{\inc_{m,n}^\ast} & C^{\D(m)\vee\D(n)}
    \end{tikzcd}
  \end{center}
  \phantom{M} % This is a hack to get the QED square to get on almost the correct line, but is there a better way?
\end{proof}

We will end this section by showing that tangential modules admit a powerful linearization trick.

\begin{theorem}\label{thm:linearize-lift}
  Suppose we have the following square of tangential $R$-modules
  \begin{center}
    \begin{tikzcd}
      A \ar{r}{u} \ar{d}{m} & X \ar{d}{e} \\
      B \ar{r}{v} & Y
    \end{tikzcd}
  \end{center}
  If there merely exists a non-linear but pointed solution $h':(B,0) \rightarrow_\ast (X,0)$ to the lifting problem above, then there merely exists a linear solution as well.
\end{theorem}

\begin{proof}
  Assume we have a non-linear solution $h':B \rightarrow X$.
  Since $\can$ is a natural isomorphism on the given square, it is sufficient to find such a linear lift on the tangent spaces.
  By \Cref{prp:derivative-inf-lin}, $\der{h'}{0}$ is the desired linear map.
\end{proof}

\begin{remark}
  Explicitly, the linear solution that the above proof produces is $h \coloneq \can^{-1} \circ \der{h'}{0} \circ \can$.
\end{remark}


\section{Adequate Modules}
A finitely presented module is described by finitely many generators and relations.
These modules can be investigated by methods of linear algebra and play a central role in algebraic geometry.
When the base ring is coherent, this notion coincides with that of coherent modules, which is usually the generality one restricts to classically.

Synthetically, however, our base ring $R$ is provably not coherent~\cite[3.2]{topology-draft} and these two notions diverge.
As it turns out, neither is well-behaved in our setting.
The category of finitely presented $R$-modules lacks kernels, and the category of coherent $R$-modules, while always abelian, is much too small as it does not even contain crucial modules like $R^n$.
In this section, we study a larger category of finitely adequate modules~(\Cref{def:fin-adeq}) as a suitable replacement whose objects still retain a certain sense of finiteness.
This notion is not completely novel and has already been studied externally in~\cite[Chapter 06Z1]{stacks-project}.

The work in this section profited from discussions with Hugo Moeneclaey, Mohamed Barakat, and Marc Nieper-Wißkirchen.

\subsection{Finitely Adequate Modules Form an Abelian Category}

In this subsection, we will show that the category of finitely adequate $R$-modules, denoted by $\finadeq$, is abelian.
We begin by recounting the notion of quasi-coherence, introduced by Ingo Blechschmid~\cite{ingo-thesis}.
For this, recall that for an $R$-module $M$ and a finitely presented $R$-algebra $A$, there is an evaluation map:
\begin{align*}
  \ev : M \otimes A & \rightarrow M^{\Spec A} \\
  m \otimes a & \mapsto (f \mapsto f(a)m)
\end{align*}

\begin{definition}\label{def:quasi-coherent}
  Let $M$ be an $R$-module.
  $M$ is \notion{quasi-coherent (qc)} if for every finitely presented $R$-algebra $A$, the map $\ev$ is an equivalence.
  We say that $M$ is \notion{weakly quasi-coherent (wqc)} if this map is an equivalence for algebras of the form $R[r^{-1}]$, for each $r : R$.
\end{definition}

We can show that quasi-coherent $R$-modules $Q$ are tangential by a direct computation.

% \rednote{These pointed maps are already linear. Does linearity make sense when the domain is $\D(n)$? Respects scalar multiplication and addition whenever it exists.}

\begin{lemma}\label{lem:disk-qc}
  Let $Q$ be a quasi-coherent $R$-module.
  A pointed map $p : (Q,0)^{\D(n)}$ is an $n$-variate linear monomial with coefficients in $Q$.
\end{lemma}

\begin{proof}
  By quasi-coherence, we have the following equivalence $Q^{\D(n)} \simeq Q \otimes \quot{R[x_1,\dots,x_n]}{\ideal{x_ix_j}} \simeq Q^{n+1}$.
  We can thus identify the maps $\D(n) \rightarrow Q$ with $n$-variate polynomials with coefficients in $Q$ and which are of degree at most $1$.
  Let $p(x_1, \dots, x_n)$ be such a polynomial; then the canonical evaluation map acts like $\ev (p(x_1, \dots, x_n))(\varepsilon_1, \dots, \varepsilon_n) = p(\varepsilon_1, \dots, \varepsilon_n)$, where $\varepsilon_i : \D (n)$.
  Whenever $p$ represents a pointed map, we have $p(0,\dots,0) = p_0 = 0$.
\end{proof}

\begin{proposition}\label{prp:tangential-qc}
  Let $Q$ be a quasi-coherent $R$-module.
  Then $Q$ is tangential.
\end{proposition}

\begin{proof}
  By the above lemma, \Cref{lem:disk-qc}, we identify $\tangent{Q}{0}$ with univariate polynomials in $Q$ of degree $1$ and constant term $0$.
  Let $q : Q$ be a point in $Q$.
  Then $\can(q)(\varepsilon) = q\varepsilon$, where $\varepsilon : \D(1)$.
  We deduce that $\can(q) = qx$, so $\can$ is an equivalence by assumption.
\end{proof}

\begin{corollary}\label{cor:inf-lin-qc}
  Let $Q$ be a quasi-coherent $R$-module.
  Then $Q$ is infinitesimally linear at $0$.
\end{corollary}

\begin{proof}
  This follows immediately from \Cref{prp:tangential-qc} and \Cref{cor:tangential-inf-lin}.
\end{proof}

Since quasi-coherent modules are tangential, they form the basis of a well-behaved category of modules.
As tangentiality is preserved by finite limits and finite colimits, we will be especially interested in the kernel closure of this category.

\begin{definition}\label{def:adequate}
  An $R$-module $M$ is \notion{adequate} if it is the kernel of a linear map $f : P \rightarrow Q$ where $P$ and $Q$ are quasi-coherent $R$-modules.
\end{definition}

Adequate $R$-modules enjoy the following properties as they form the kernel closure of quasi-coherent modules.

\begin{proposition}\label{prp:wqc-adeq}
  Adequate $R$-modules are weakly quasi-coherent. \earlyqed
\end{proposition}

\begin{proposition}\label{prp:adeq-tangential}
  Adequate $R$-modules are tangential. \earlyqed
\end{proposition}

For the remainder of this section, we restrict our attention to a subcategory of adequate modules that satisfy certain finiteness conditions. We will construct the objects of this subcategory by first introducing their basic building blocks:

\begin{definition}\label{def:lin-adeq}
  Let $L$ be an $R$-module.
  $L$ is called \notion{linearly adequate} if it is the kernel of a linear map $f : R^m \rightarrow R^n$.
  We refer to the map $f$ as a \notion{linear copresentation} of $L$.
\end{definition}

\begin{remark}\label{rem:lin-adeq}
  In the article \cite{diffgeo-article}, linearly adequate modules went under the name finitely copresented modules.
  We decided to change the name here as linearly adequate matches the name of the external concept \cite[Chapter 06Z1]{stacks-project}.
\end{remark}

\begin{definition}\label{def:dual}
  Let $M$ be an $R$-module.
  We define the \notion{dual} of $M$ to be the module $M^\ast \defeq \Hom_R(M,R)$.
\end{definition}

\begin{theorem}[{{\cite[Lemma 2.3.9, Corollary 2.3.10]{diffgeo-article}}}]\label{thm:lin-adeq-dual}
  Dualization restricts to an anti-equivalence between linearly adequate modules and finitely presented ones.
  In other words, dualizing a linear copresentation of $L$,
  \[ M \hookrightarrow R^n \to R^m \]
  yields a finite presentation of $L^\ast$
  \[ R^m \to R^n \twoheadrightarrow L^\ast \]
  Moreover, the canonical map into the double dual, $\ev : L \to L^{\ast\ast}$, is an isomorphism for finitely presented and linearly adequate $R$-modules. \earlyqed
\end{theorem}

\begin{lemma}\label{lem:fp-lift}
  Let $M$ and $N$ be finitely presented $R$-modules and $f: M \to N$ a linear map.
  Then, there is an extension of this morphism to all presentations of $M$ and $N$:
  \begin{center}
    \begin{tikzcd}
      R^{m'}\ar[r]\ar[d,dotted,"\exists"] & R^m\ar[->>,r,"\pi_M"]\ar[d,dotted,"\exists"] & M\ar[d, "f"] \\
      R^{n'}\ar[r] & R^n\ar[->>,r,"\pi_N"] & N \\
    \end{tikzcd}
  \end{center}
  By duality, the analogous statement holds for linearly adequate modules.
\end{lemma}

\begin{proof}
  This is immediate since finite free $R$-modules are projective in $\modcat{R}$.
\end{proof}

From every commutative square, we can extract some exact sequences.
These will be particularly useful to show that finitely presented $R$-modules are closed under cokernels and linearly adequate $R$-modules are closed under kernels.

\begin{lemma}\label{lem:diff-ker-is-ker}
  Suppose we have a commutative square:
  \begin{center}
    \begin{tikzcd}
      A \ar{r}{u} \ar{d}{f} & B \ar{d}{g} \\
      C \ar{r}{v} & D
    \end{tikzcd}
  \end{center}
  The kernel $K$ of the induced map $\Ker(u) \rightarrow \Ker(v)$ is equivalent to the kernel $K'$ of $(u, f) : A \to B \oplus C$.
\end{lemma}

\begin{proof}
  Observe that $\Ker(u) \subseteq \Ker(gu)$.
  Thus $\Ker(u) = \Ker((u, gu) : A \to B \oplus D)$.
  Therefore, we obtain this result by the third isomorphism theorem:
  \begin{center}
    \begin{tikzcd}
      K \ar{r}{\simeq} \ar[hook]{d} & K' \ar[hook]{d} \\
      \Ker(u) \ar{d} \ar[hook]{r} & A \ar{d}{(u,f)} \ar{r}{(u,gu)} & B \oplus D \ar[equal]{d} \\
      \Ker(v) \ar[hook]{r} & B \oplus C \ar{r}{\mathrm{id} \oplus v} & B \oplus D
    \end{tikzcd}
  \end{center}
  \hfill{\phantom{.}}
\end{proof}

\begin{proposition}\label{prp:lin-adeq-ker-closure}
Linearly adequate $R$-modules are closed under kernels.
Dually, finitely presented $R$-modules are closed under cokernels.
\end{proposition}

\begin{proof}
  Suppose that $K$ is the kernel of a linear map $f : M \rightarrow N$, where $M$ and $N$ are linearly adequate.
  By the dual of \Cref{lem:fp-lift}, we can find a linear extension of $f$ to copresentations of $M$ and $N$:
  \begin{center}
    \begin{tikzcd}
      K \ar[hook]{d} \\
      M \ar{d} \ar[hook]{r} & R^{m_1} \ar{r} \ar{d} & R^{m_2} \ar{d} \\
      N \ar[hook]{r} & R^{n_1} \ar{r} & R^{n_2}
    \end{tikzcd}
  \end{center}
  By \Cref{lem:diff-ker-is-ker}, $K$ is also the kernel of a map $R^{m_1} \rightarrow R^{m_2 + n_1}$, so it is itself linearly adequate.
\end{proof}

\begin{lemma}\label{lem:ker-fp-is-cok-lin-adeq}
  If $f : M \to N$ is a map of finitely presented $R$-modules, then $\ker f$ is the cokernel of a map between linearly adequate $R$-modules.
\end{lemma}

\begin{proof}
  By \Cref{lem:fp-lift}, we can assume that $f$ is induced by a square of morphisms between finite free $R$-modules and construct the linearly adequate modules $A$ and $B$ as described below:
  \begin{center}
  \begin{tikzcd}
    A\ar[r,"\pi_1"]\ar[d,"l\times\id"] & R^{m'}\ar[r,"t"]\ar[d,"l"] & R^{n'}\ar[d,"r"]   \\
    B\ar[r,"\pi_1"]\ar[d] & R^m\ar[r,"b"]\ar[d] & R^n\ar[d] \\
    K\ar[r] &M\ar[r,"f"] & N
    \end{tikzcd}
   \end{center}
   $A:=\{(z,y):R^{m'}\times R^{n'}\mid blz=ry\}$, $B:=\{(x,y):R^{m}\times R^{n'}\mid bx=ry\}$.
\end{proof}

% We name the following modules after Adelman, since it will turn out that what we construct agrees with the free abelian category on the finite free $R$-modules, and our construction is close to Adelman's general construction of the free abelian category.

The modules in the kernel-cokernel closure of the finite free $R$-modules might yield a replacement of the algebro-geometric notion of coherent sheaf of modules, which would be a trivial notion over $R$.

\begin{definition}\label{def:fin-adeq}
  Let $M$ be an $R$-module. Then $M$ is a \notion{finitely adequate $R$-module} if one of the following equivalent statements holds:
  \begin{enumerate}[(i), noitemsep]
    \item $M$ is merely the kernel of a homomorphism between finitely presented $R$-modules.
    \item $M$ is merely the cokernel of a homomorphism between linearly adequate $R$-modules.
  \end{enumerate}
  We denote the type of finitely adequate modules with $\finadeq$\index{$\finadeq$}, as well as the category of these modules.
\end{definition}

\begin{proposition}\label{prp:fin-adeq-adeq}
  Every finitely adequate $R$-module is adequate.
  In particular, finitely adequate $R$-modules are also tangential and wqc.
\end{proposition}

\begin{proof}
  Every finitely presented $R$-module is quasi-coherent.
  Since every finitely adequate $R$-module is merely the kernel of a map between quasi-coherent modules, it is itself adequate.
  The rest follows from \Cref{prp:tangential-qc} and \Cref{prp:wqc-adeq}.
\end{proof}

We will first characterize the projective and injective objects in $\finadeq$.

\begin{lemma}\label{lem:lin-adeq-lift}
  Let $L$ be a linearly adequate $R$-module.
  Every lifting problem
  \begin{center}
    \begin{tikzcd}
      & A \ar[two heads]{d}{e} \\
      L \ar{r}{v} & B
    \end{tikzcd}
  \end{center}
  merely has a linear solution when $A$ and $B$ are tangential and $M \defeq \Ker e$ is wqc.
\end{lemma}

\begin{proof}
  By \cite[Theorem 7.3.6]{draft} we have $H^1(L,M)=0$, since $L$, like any linearly adequate $R$-module, is an affine scheme. Indeed, as the kernel of a linear map $f : R^m \to R^n$, the module $L$ is equivalent to $\Spec(\quot{R[x_1, \dots, x_m]}{\langle f(x_1), \dots, f(x_n) \rangle})$.

  By definition, we have $H^1(L,M)=||L \to K(M,1)||_0$ and therefore $L \to K(M,1)$ is connected.
  For any $b : B$, we let $E_b$ denote the fiber of $e$ over $b$.
  Then $E_b$ is an $M$-torsor, or in other words an element of $K(M,1)$.
  So $E \circ v$ defines a map $L \to K(M,1)$, and we have $||E \circ v = (\argmhole \mapsto M)||$ by the connectedness noted above.
  This means that the dependent type $E \circ v$ merely has a section, since $(\argmhole \mapsto M)$ has a section.
  Let $h'' : L \rightarrow A$ denote the section after projecting onto $A$.
  The map $h' \defeq h'' - h''(0)$ is now a pointed map of type $(L,0) \rightarrow_\ast (A,0)$.
  The map $h'$ still solves the problem as $e(h''(0)) = 0$.
  By \Cref{thm:linearize-lift}, there is a linear map $h : L \rightarrow A$ solving this lifting problem.
\end{proof}

\begin{theorem}\label{thm:lin-adeq-proj}
  Let $M$ be a finitely adequate module.
  Then $M$ is projective in $\finadeq$ if and only if $M$ is linearly adequate.
\end{theorem}

\begin{proof}
  Assume first that $M$ is linearly adequate.
  We need to merely solve any lifting problem of the form
  \begin{center}
    \begin{tikzcd}
      & A \ar[two heads]{d}{e} \\
      M \ar{r}{v} & B
    \end{tikzcd}
  \end{center}
  where $A$ and $B$ are finitely adequate.
  Such a solution follows from \Cref{lem:lin-adeq-lift}, since $A$ and $B$ are tangential and wqc by \Cref{prp:fin-adeq-adeq}.

  Assume now that $M$ is projective.
  Let $L_1 \xrightarrow{p} L_2 \twoheadrightarrow M$ be a presentation of $M$ by linearly adequate modules, and let $M' \defeq \Im p$.
  Since $M$ is assumed projective, the epimorphism $L_2 \twoheadrightarrow M$ splits.
  It follows that the inclusion $M' \hookrightarrow L_2$ splits as well and $M'$ is a direct summand of $L_2$.
  Since $L_2$ is linearly adequate, it is projective, thus $M'$ is projective as well.
  Since $M'$ is projective, the canonical map $L_1 \twoheadrightarrow M'$ splits.
  This gives us a map $L_2 \rightarrow L_1$ such that $M$ is its kernel.
  $M$ is now linearly adequate by \Cref{prp:lin-adeq-ker-closure}, as it is the kernel of a map between linearly adequate $R$-modules.
\end{proof}

\begin{lemma}\label{lem:fin-adeq-lift}
  Any linear map $f : M \rightarrow N$ merely lifts to a commutative square of linearly adequate modules.
\end{lemma}

\begin{proof}
  This is a direct consequence of \Cref{thm:lin-adeq-proj}.
\end{proof}

\begin{proposition}\label{prp:fin-adeq-cok-closure}
  Finitely adequate modules are closed under cokernels.
\end{proposition}

\begin{proof}
  This proof follows from \Cref{lem:diff-ker-is-ker} and \Cref{lem:fin-adeq-lift}, following a proof analogous to \Cref{prp:lin-adeq-ker-closure}.
\end{proof}

\begin{lemma}\label{lem:dual-exact-fp-inc}
  Let $M$ be a finitely adequate $R$-module, and $P$ a finitely presented $R$-module.
  Any lifting problem
  \begin{center}
    \begin{tikzcd}
      M \ar[hook]{d}{m} \ar{r}{u} & R \\
      P
    \end{tikzcd}
  \end{center}
  merely admits a linear solution.
\end{lemma}

\begin{proof}
  By \Cref{prp:fin-adeq-cok-closure}, the cokernel $N \defeq \Cok(m)$ is finitely adequate.
  Thus, there exists a monomorphism $n : N \hookrightarrow Q$ such that $Q$ is finitely presented.
  Let $f : P \rightarrow Q$ be the map induced by $n$.
  Now, $m$ is the kernel of $f$.
  By \Cref{lem:ker-fp-is-cok-lin-adeq}, we obtain an analysis of $M$:
  \begin{center}
    \begin{tikzcd}
      K \ar{d}{p \times \id} \ar[hook]{r} & R^{m_1} \ar{r}{f_1} \ar{d}{p} & R^{n_1} \ar{d}{q} \\
      L \ar[two heads]{d} \ar[hook]{r} & R^{m_2} \ar{r}{f_2} \ar[two heads]{d} & R^{n_2} \ar[two heads]{d} \\
      M \ar[hook]{r}{m} & P \ar{r}{f} & Q
    \end{tikzcd}
  \end{center}
  In the diagram above, we have $P = \Cok(p)$, $Q = \Cok(q)$, $K = R^{m_1} \underset{R^{n_2}}{\times} R^{n_1}$, and $L = R^{m_2} \underset{R^{n_2}}{\times} R^{n_1}$.
  Let $t$ be the composite $f_2 \circ p = q \circ f_1$.
  
  In the above analysis, only the columns are exact.
  We will alter this diagram so that every column and row becomes exact.
  \begin{center}
    \begin{tikzcd}
      K \ar{d}{p \times \id} \ar[hook]{r} & R^{m_1 + n_1} \ar{r}{t + q} \ar{d}{p \times \id} & R^{n_2} \ar[equal]{d} \\
      L \ar[two heads]{d} \ar[hook]{r} & R^{m_2 + n_1} \ar{r}{f_2 + q} \ar[two heads]{d} & R^{n_2} \\
      M \ar[hook]{r}{m} & P
    \end{tikzcd}
  \end{center}

  The linear map $u : M \to R$ restricts to a map $l_1 : L \to R$ that is zero on $K$.
  By~\cite[Lemma 2.3.7]{diffgeo-article}, it is possible to extend $l_1$ to $l_2 : R^{m_2+n_1} \to R$.
  $l_2$ will descend to $M$ if and only if it vanishes on the image of $p \times \id$.
  Although this is not generally the case, we can construct a linear correction $c : R^{m_2} \to R$ such that $l_2 - c$ has this property.

  To construct $c$, let $c_0 \coloneq l_2 \circ (p \times \id)$.
  Note that $c_0$ vanishes on $K$ and therefore descends to a linear map $c_1 : \Im(t + q) \to R$.
  By \cite[Lemma 2.3.8]{diffgeo-article}, $c_1$ extends to a map $c_2 : R^{n_2} \to R$.
  Now, $c \coloneq c_2 \circ (f_2 + q)$ is a linear map on $R^{m_2 + n_1}$ that is zero on $K$,
  so $l_2 - c$ vanishes on the image of $p \times \id$.
  Therefore, $l_2 - c$ descends to a map $l : P \to R$.

  Finally, note that $l$ restricts to $l_1$ on $L$.
  Thus we have $u = l \circ m$.
\end{proof}

Let $M$ be an $R$-module.
There is a natural map $\ev : M \rightarrow M^{\ast\ast}$ that sends a term $m$ to $\ev_m(f) \defeq f(m)$.

\begin{theorem}\label{thm:can-dual-iso}
  Let $M$ be a finitely adequate module.
  Then $\ev : M \rightarrow M^{\ast\ast}$ is an equivalence.
\end{theorem}

\begin{proof}
  We will prove this by exploiting two different naturality squares involving $\ev$.

  Since $M$ is finitely adequate, it merely embeds into a finitely presented $R$-module $P$.
  This yields the naturality square
  \begin{center}
    \begin{tikzcd}
      M \ar[hook]{r} \ar{d}[right]{\ev} \ar[hook]{rd} & P \ar{d}[right]{\ev}[left]{\simeq} \\
      M^{\ast\ast} \ar{r} & P^{\ast\ast}
    \end{tikzcd}
  \end{center}
  Since the composite $M \rightarrow P^{\ast\ast}$ is a monomorphism, so $\ev : M \rightarrow M^{\ast\ast}$ must be a monomorphism as well.

  Since $M$ is finitely adequate, it is merely covered by a linearly adequate $R$-module $L$.
  Observe that the dual of the cover, $M^\ast \hookrightarrow L^\ast$, is a monomorphism.
  Taking the dual again yields an epimorphism $L^{\ast\ast} \twoheadrightarrow M^{\ast\ast}$, since $L^{\ast}$ is finitely presented by \Cref{thm:lin-adeq-dual} and \Cref{lem:dual-exact-fp-inc}.
  This yields the naturality square
  \begin{center}
    \begin{tikzcd}
      L \ar[two heads]{r} \ar{d}[right]{\ev}[left]{\simeq} \ar[two heads]{rd} & M \ar{d}{\ev} \\
      L^{\ast\ast} \ar[two heads]{r} & M^{\ast\ast}
    \end{tikzcd}
  \end{center}
  Since the composite $L \rightarrow M^{\ast\ast}$ is an epimorphism, $\ev : M \rightarrow M^{\ast\ast}$ has to be an epimorphism as well.
\end{proof}

\begin{corollary}\label{cor:dual-equiv}
  Taking duals $(\argmhole)^{\ast} : \finadeq \rightarrow \finadeq$ is an anti-equivalence of categories. \earlyqed
\end{corollary}

We get the following results following duality.

\begin{theorem}\label{thm:fp-inj}
  Let $M$ be finitely adequate.
  Then $M$ is injective in $\finadeq$ if and only if $M$ is finitely presented.
\end{theorem}

\begin{proof}
  By \Cref{cor:dual-equiv} it is sufficient to prove the dual statement: $M$ is projective in $\finadeq$ if and only if $M$ is linearly adequate.
  This is precisely \Cref{thm:lin-adeq-proj}.
\end{proof}

\begin{corollary}\label{cor:fin-adeq-extension}
  Any linear map $f : M \rightarrow N$ merely extends to a commutative square of finitely presented $R$-modules. \earlyqed
\end{corollary}

\begin{corollary}\label{cor:fin-adeq-ker-closure}
  Finitely adequate modules are closed under kernels. \earlyqed
\end{corollary}

Since the category of finitely adequate modules is a subcategory of $\modcat{R}$ which is closed under finite limits and finite colimits, it is itself abelian.

\begin{theorem}\label{thm:fin-adeq-abelian}
  The category $\finadeq$ is abelian. \earlyqed
\end{theorem}

\begin{theorem}\label{thm:fin-adeq-class}
  Let $M$ be an $R$-module.
  The following are equivalent:
  \begin{enumerate}[(i), noitemsep]
    \item $M$ is finitely adequate.
    \item $M$ is the image of a map from a linearly adequate $R$-module to a finitely presented $R$-module.
    \item $M$ is in the smallest abelian subcategory of $\modcat{R}$ that contains all finite free $R$-modules.
  \end{enumerate}
\end{theorem}

\begin{proof}
  The equivalence between (i) and (ii) is clear from \Cref{prp:fin-adeq-cok-closure} and \Cref{cor:fin-adeq-ker-closure}.
  The equivalence with (iii) will follow from \Cref{thm:fin-adeq-abelian}.
\end{proof}

\begin{remark}\label{rem:fin-adeq-class-self-dual}
  Observe that condition (ii) is self-dual.
\end{remark}

\begin{remark}\label{rem:fin-adeq-Adelman}
  By \cite[Theorem 1.15]{adelman-construction}, \Cref{thm:lin-adeq-proj} and \Cref{thm:fp-inj} imply that the category $\finadeq$ is the free abelian category over the finite free $R$-modules.
  This means the finitely adequate $R$-modules have the following universal property:
  For every abelian category $\mathcal{A}$ and all additive functors $F:\finfree \to \mathcal{A}$, there exists a unique exact functor $K$ as follows:
  \begin{center}
    \begin{tikzcd}
      \finfree \ar{d} \ar{r}{F} & \mathcal{A} \\
      \finadeq \ar[dashed]{ru}[swap]{K}
    \end{tikzcd}
  \end{center}
  Consequently, the data needed to specify an exact functor $K$ is an assignment $KR$ of $R$ in $\mathcal{A}$ and a group homomorphism $\Hom_R(R,R) \rightarrow \mathcal{A}(KR,KR)$.
\end{remark}

\subsection{Properties of Finitely Adequate Modules}

In this subsection we collect some properties of finitely adequate $R$-modules.

\begin{proposition}
  Let $A$ be a closed submodule of a finitely adequate module $M$. Then $A$ is finitely adequate if $A$ is tangential.
\end{proposition}

\begin{proof}
  Let $M$ be presented by the linearly adequate modules $K \to L \twoheadrightarrow M$. Consider the pullback $L' \coloneq A \times_M L$.
  Since pullbacks of $R$-modules are taken as pullbacks on the underlying sets, we have that $L'$ is a closed submodule of $L$, and is therefore a scheme by \cite[Proposition 5.4.3]{draft}.
  Therefore the tangent space $\tangent{L'}{0}$ is a linearly adequate module by \cite[Lemma 2.2.8]{diffgeo-article}.
  Since tangential modules are closed under limits \Cref{prp:op-closure-inf-lin-tangential}, $L'$ is also tangential, and therefore also linearly adequate.
  The sequence $K \to L' \twoheadrightarrow A$ presents $A$ as a cokernel of a map of linearly adequate modules, so $A$ is finitely adequate.
\end{proof}

\begin{proposition}\label{prp:fin-adeq-hom-closure}
  For finitely adequate $R$-modules $M$ and $N$, $\Hom_R(M,N)$ is a finitely adequate $R$-module as well.
\end{proposition}

\begin{proof}
  The proof in \cite[Chapter IV, 4.12]{lombardi-quitte} works, if ``coherent'' is replaced by ``finitely adequate''.

  First we use their proof to show that for finitely presented $R$-modules $M$ and $N$, the $R$-module $\Hom_R(M,N)$ is finitely adequate. Morphisms $f: M \to N$ are presented by squares:
  \begin{center}
    \begin{tikzcd}
      R^m \ar{r} \ar{d}{\varphi} & R^{m'} \ar{d}{\psi} \\
      R^n\ar{r} & R^{n'}
    \end{tikzcd}
  \end{center}
  So we have a finite free module of pairs of morphisms $(\varphi,\psi)$ and a submodule of pairs such that the square above commutes.
  This submodule $S$ is the kernel of a linear map and therefore finitely adequate.
  We have a surjection $\pi : S \to \Hom_R(M,N)$, so $\Hom_R(M,N)$ is finitely adequate if $\pi$ is a cokernel of a map of finitely adequate $R$-modules.
  This is the case, since there is a surjection onto $\Ker{\pi}$ from the finite free $R$-module of linear maps which split the square:
  \begin{center}
    \begin{tikzcd}
      R^m \ar{r} \ar{d}{\varphi} & R^{m'} \ar{d}{\psi} \ar{ld}{s} \\
      R^n \ar{r} & R^{n'}
    \end{tikzcd}
  \end{center}
  Then we can reuse the same argument to show the statement of the proposition.
\end{proof}

\begin{proposition}\label{prop:hood-faithful}
  The functor $\hood{\argmhole}{1} : \modcat R \rightarrow \pType$ is faithful when restricted to the subcategory of tangential $R$-modules.
\end{proposition}

\begin{proof}
  Let $f,g : M \rightrightarrows N$ be two linear maps between tangential $R$-modules with $\hood{f,0}{1} = \hood{g,0}{1}$.
  Then $df_0=dg_0$ and by naturality of the map $\can$ the claim follows.
  % This is immediate, as the action on the map is given by restriction, i.e. $\hood{f,0}{1} \defeq f\mid_{\hood{\argmhole}{1}}$.
  % To see that it is full, let $M$ and $N$ be finitely adequate $R$-modules, and let $f : \hood{M,0}{1} \rightarrow_\ast \hood{N,0}{1}$ be a pointed map.
  % Because $M$ and $N$ are finitely adequate, we can find linearly adequate $R$-modules $K_1$, $K_2$, $L_1$, and $L_2$ together with presentations $K_1 \rightarrow K_2 \twoheadrightarrow M$ and $L_1 \rightarrow L_2 \twoheadrightarrow N$.
  % By ???, $\hood{K_i,0}{1}$ and $\hood{L_i,0}{1}$ are projective types.
  % Using this property, we can lift the map $f$ to a commutative square by \Cref{lem:hood-module-prs-surj}
  % \begin{center}
  %   \begin{tikzcd}
  %     \hood{K_1,0}{1} \ar{d} \ar[dashed]{r} & \hood{L_1,0}{1} \ar{d} \\
  %     \hood{K_2,0}{1} \ar[two heads]{d} \ar[dashed]{r} & \hood{L_2,0}{1} \ar[two heads]{d} \\
  %     \hood{M,0}{1} \ar{r}{f} & \hood{N,0}{1}
  %   \end{tikzcd}
  % \end{center}
  % By \cite[Lemma 2.1.4]{diffgeo-article} and \Cref{thm:lin-adeq-dual}, the morphisms in the top square extend to modules.
  % By taking cokernels we get an induced map
  % \begin{center}
  %   \begin{tikzcd}
  %     K_1 \ar{r} \ar{d} & L_1 \ar{d} \\
  %     K_2 \ar[two heads]{d} \ar{r} & L_2 \ar[two heads]{d} \\
  %     M \ar{r}{f'} & N
  %   \end{tikzcd}
  % \end{center}
  % It is clear that $f'\mid_{\hood{M,0}{1}} = f$.
\end{proof}

\begin{remark}
  The functor $\hood{\argmhole}{1}$ of the previous proposition is full when restricted even further to linearly adequate $R$-modules \rednote{cite}.
  This, however, already fails on finitely adequate $R$-modules.

  To see this, consider the maps $\alpha_\varepsilon : (\varepsilon) \to (\varepsilon), r\varepsilon \mapsto r^2\varepsilon$ for every $\varepsilon:\D(1)$.
  These are well-defined, since $r\varepsilon = s\varepsilon$ implies $r^2\varepsilon-s^2\varepsilon = (r+s)(r-s)\varepsilon = 0$.
  Since for every $n:(\varepsilon)$ we have $(\varepsilon=0)\to(n=0)$, we conclude that $(\varepsilon) = \hood{(\varepsilon),0}{1}$.
  Note that $\hood{\argmhole}{1}$ is the identity functor on those $R$-modules that are their own first order neighbourhood.
  Thus it suffices to show that it is not the case that all $\alpha_\varepsilon$ are $R$-linear.
  Assume they are.
  Then $r^2\varepsilon = \alpha_\varepsilon(r\varepsilon) = r\alpha_\varepsilon(\varepsilon) = r\varepsilon$ and hence also $(r^2-r)\varepsilon=0$ for all $r:R$.
  Using duality and comparing coefficients, we see that $\varepsilon=0$ for every $\varepsilon:\D(1)$, which cannot be.
\end{remark}

\rednote{The next remark could maybe move to the appendix.  Not sure if we need it in the final paper.}

\begin{remark}\label{rem:hood-module-prs-surj}
  If $f : M \rightarrow N$ be an epimorphism of $R$-modules with $M$ linearly adequate
  and $N$ finitely adequate,
  it does not follow that $\hood{f,0}{1}$ is surjective.
  Here is an example.

  Let $\varepsilon : \D(1)$ be arbitrary.
  Consider the map $\pi_\varepsilon : R \to (\varepsilon)$ which sends $r$ to $r \varepsilon$.
  Note that $R$ is linearly adequate and $(\varepsilon)$ is finitely adequate
  by \Cref{thm:fin-adeq-class}.
  We will show that $\pi_\varepsilon$ does not satisfy the claim for all $\varepsilon$.
  Indeed, assume that $\hood{\pi_{\varepsilon}, 0}{1}$ is surjective for all $\varepsilon$.
  Note that $\varepsilon \in \hood{(\varepsilon), 0}{1} = (\varepsilon)$.
  Let $r:\D(1) = \hood{R,0}{1}$ be a lift of $\varepsilon$.
  Since $\pi_\varepsilon(1) = \varepsilon$, $\pi_\varepsilon(r-1)=0$, i.e., $(r-1)\varepsilon=0$.
  But $r$ is nilpotent, hence $r-1$ invertible and this would mean $\varepsilon = 0$, which cannot be true for all $\varepsilon$.
\end{remark}

Finitely adequate modules form a weak Serre subcategory of $\modcat{R}$.

\begin{lemma}\label{lem:fin-adeq-ext-pb-red}
  Suppose we have a commutative diagram
  \begin{center}
    \begin{tikzcd}
      B \ar[equal]{d} \ar[hook]{r} & p^*E \ar[two heads]{r} \ar[two heads]{d} \ar[phantom]{rd}[very near start]{\lrcorner}[very near end]{\ulcorner} & P \ar[two heads]{d}{p} \\
      B \ar[hook]{r} & E \ar[two heads]{r} & A
    \end{tikzcd}
  \end{center}
  in which $B$, $A$ and $P$ are finitely adequate $R$-modules.
  Then $E$ is finitely adequate if and only if $p^*E$ is so.
\end{lemma}

\begin{proof}
  Suppose that $E$ is finitely adequate.
  Since $p^*E$ is a pullback of finitely adequate $R$-modules, it is itself finitely adequate.

  Now assume that $p^*E$ is finitely adequate.
  By the third isomorphism theorem we can prove that $E$ is finitely adequate.
  In the following diagram, the object $0^*_Ap^*E$ is the kernel of $p^*E \rightarrow E$, $0^*_AE$ is the kernel of $E \rightarrow A$, and $0^*_Ep^*E$ is the kernel of $p^*E \rightarrow E$.
  \begin{center}
    \begin{tikzcd}
      0^*_Ep^*E \ar[equal]{r} \ar[hook]{d} & 0^*_Ep^*E \ar[hook]{d} \\
      0^*_Ap^*E \ar[hook]{r} \ar[two heads]{d} & p^*E \ar[two heads]{r} \ar[two heads]{d} & A \ar[equal]{d} \\
      B \ar[hook]{r} & E \ar[two heads]{r} & A
    \end{tikzcd}
  \end{center}
  Chasing around this diagram we can see that $0^*_Ap^*E$ is finitely adequate, because $p^*E$ is.
  Furthermore, $0^*_Ep^*E$ is finitely adequate, because $0^*_Ap^*E$ is finitely adequate.
  It then follows that $E$ is finitely adequate, because it is the cokernel of two finitely adequate $R$-modules.
\end{proof}

\begin{lemma}\label{lem:fin-adeq-ext-red}
  Finitely adequate $R$-modules are closed under extension if and only if for any linearly adequate $R$-module $L$ and finitely adequate $R$-module $B$, any exact sequence $0 \to B \to E \to L \to 0$ splits.
\end{lemma}

\begin{proof}
  Assume that finitely adequate modules are closed under extensions.
  Then it is clear that any extension by $B$ and $L$ is split, as $L$ is injective among the finitely adequate modules \Cref{thm:lin-adeq-proj}.

  Assume now that $\textup{Ext}^1(L,B) = 0$ for any such modules $L$ and $B$.
  Consider an extension problem
  \begin{center}
    \begin{tikzcd}
      B \ar[hook]{r} & E \ar[two heads]{r} & A
    \end{tikzcd}
  \end{center}
  where $B$ is finitely adequate, and $L$ is finitely adequate and is merely presented by $L_1 \rightarrow L_2 \overset{p}{\twoheadrightarrow} A$.
  By \Cref{lem:fin-adeq-ext-pb-red}, we can take the pullback of the extension along $p : P_1 \twoheadrightarrow L$ to get another associated extension.
  \begin{center}
    \begin{tikzcd}
      B \ar[equal]{d} \ar[hook]{r} & p^*E \ar[two heads]{d} \ar[two heads]{r} \ar[phantom]{rd}[very near start]{\lrcorner} & L_1 \ar[two heads]{d} \\
      B \ar[hook]{r} & E \ar[two heads]{r} & A
    \end{tikzcd}
  \end{center}
  Now $E$ is finitely adequate if and only if $p^*E$ is so.

  Since we assumed that any exact sequence on the form $0 \to B \to E \to L_1 \to 0$ splits, it is necessary that $p^*B = B \oplus L_1$, which is finitely adequate.
  Thus, $E$ is finitely adequate as desired.
\end{proof}

\begin{theorem}\label{thm:fin-adeq-ext-closed}
  Finitely adequate $R$-modules are closed under extension.
\end{theorem}

\begin{proof}
  By \Cref{lem:fin-adeq-ext-red} it is sufficient to show that any exact sequence on the form $0 \to B \to E \to L \to 0$ splits for any finitely adequate $B$ and linearly adequate $L$.
  This amounts to showing that any lifting problem
  \begin{center}
    \begin{tikzcd}
      & E \ar[two heads]{d}{p} \\
      L \ar[equal]{r} & L
    \end{tikzcd}
  \end{center}
  where $L$ is finitely adequate and $B \defeq \Ker(p)$ is finitely adequate, merely has a linear solution.
  Observe that $E$ is tangential by \Cref{prp:op-closure-inf-lin-tangential}, since $B$ and $L$ are tangential by \Cref{prp:adeq-tangential}.
  $B$ is also wqc by \Cref{prp:wqc-adeq}.
  Thus the lifting problem has a solution by \Cref{lem:lin-adeq-lift}.
\end{proof}


By \Cref{rem:tensor-not-closed}, the tensor product $M \otimes N$ is not necessarily linearly adequate when $M$ and $N$ are linearly adequate.
This contrasts with the case where $M$ is finitely presented (or, similarly, where $N$ is finitely presented), since $M \otimes N$ can then be realized as the cokernel of a map between finite direct sums of $N$, and is therefore finitely presented.
Somewhat surprisingly, the internal hom still has a left adjoint in $\finadeq$.
We will show that the \notion{double dual tensor product} $M \dbldualtensor N \defeq (M \otimes N)^{\ast\ast}$ is this left adjoint.
By the hom-tensor adjunction, there is already a natural isomorphism $\Hom_R(A, \Hom_R(B, C)) \simeq \Hom_R(A \otimes B, C)$.
Thus, it is sufficient to prove that $\ev_{A \otimes B}^\ast : \Hom_R(A \dbldualtensor B, C) \rightarrow \Hom_R(A \otimes B, C)$ is a natural isomorphism for all finitely adequate modules $A$, $B$, and $C$.
The composite of these natural isomorphisms yields the natural isomorphism $\Hom_R(A \dbldualtensor B, C) \simeq \Hom_R(A, \Hom_R(B, C))$.
It is clear that $\ev_{A \otimes B}^\ast$ is natural in all three arguments.
Now, $\ev_{A \otimes B}^\ast$ is an isomorphism precisely when every morphism $f : A \otimes B \to C$ factors uniquely through $\ev_{A \otimes B} : A \otimes B \to A \dbldualtensor B$.

We first carry out this factorization in the general setting of monads.
Let $T$ be a monad on $\mathcal{C}$ with multiplication $\mu$ and unit $\eta$.
In our situation, $\mathcal{C}$ will be the category of $R$-modules and $T$
will be the double-dual operation.

\begin{definition}\label{def:monad-reflexive}
  Let $X$ be an object of $\mathcal{C}$. We say that $X$ is:
  \begin{itemize}
    \item \notion{$T$-torsionless} if the unit $\eta_X$ is a monomorphism,
    \item \notion{$T$-reflexive} if the unit $\eta_X$ is an isomorphism.
  \end{itemize}
\end{definition}

\begin{lemma}\label{lem:monad-app-torsionless}
  Let $X$ be an object of $\mathcal{C}$. Then $TX$ is $T$-torsionless.
\end{lemma}

\begin{proof}
  By unitality, we have $\mu_X \circ \eta_{TX} = \id_{TX}$, so $\eta_{TX}$ is a split monomorphism.
\end{proof}

\begin{definition}\label{def:rel-epi}
  Let $\mathfrak{C}$ be a class of objects in $\mathcal{C}$. A map $f : X \to Y$ is \textbf{$\mathfrak{C}$-epic} if for every $Z$ in $\mathfrak{C}$ and all morphisms $g, h : Y \to Z$ such that $g \circ f = h \circ f$, we have $g = h$.
\end{definition}

The monad $T$ is said to be \notion{well-pointed} if $T\eta_X = \eta_{TX}$ for every object $X$ of $\mathcal{C}$.
We will not discuss well-pointed monads in this global sense, but rather focus on the local notion.
For a treatment of well-pointed monads, we refer the reader to Chapter 4.2 of Borceux's book~\cite{Borceux_1994_II}.
The next result is similar to Proposition 4.2.3 of Borceux~\cite{Borceux_1994_II}, but we do not assume that $T$ is well-pointed, which forces us to give a different proof.

\begin{proposition}\label{prp:monad-well-pointedness-torsionless-epi}
  Let $X$ be an object of $\mathcal{C}$. The following are equivalent:
  \begin{enumerate}[(i), noitemsep]
    \item $TX$ is $T$-reflexive,
    \item $T\eta_X = \eta_{TX}$,
    \item $\eta_X$ is $T$-torsionless-epic.
  \end{enumerate}
\end{proposition}

\begin{proof}
  (i) $\Rightarrow$ (iii): Assume that $TX$ is $T$-reflexive, i.e., that $\eta_{TX}$ is an isomorphism.
  Then $T\eta_X$ is an isomorphism as well, since $T\eta_X$ is a section of $\mu_X$.
  Consider two morphisms $f, g : TX \to Y$, where $Y$ is $T$-torsionless and $f \circ \eta_X = g \circ \eta_X$.
  By applying $T$ to this equality, we obtain $Tf \circ T\eta_X = Tg \circ T\eta_X$.
  Thus, $Tf = Tg$ since $T\eta_X$ is an epimorphism.
  By naturality of $\eta$, we have $\eta_Y \circ f = Tf \circ \eta_{TX} = Tg \circ \eta_{TX} = \eta_Y \circ g$.
  Since $Y$ is $T$-torsionless, $\eta_Y$ is a monomorphism and thus $f = g$.
  
  (iii) $\Rightarrow$ (ii):  Assume that $\eta_X$ is $T$-torsionless-epic.
  By naturality of $\eta$, we have $\eta_{TX} \circ \eta_X = T\eta_X \circ \eta_X$. 
  Since $TX$ is $T$-torsionless by \Cref{lem:monad-app-torsionless}, we obtain $T\eta_X = \eta_{TX}$.
  
  (ii) $\Rightarrow$ (i): Assume that $T\eta_X = \eta_{TX}$.
  Our goal is to show that $\eta_{TX}$ is an isomorphism, so it suffices to show that $T\eta_X$ is an isomorphism.
  By unitality, we have that $\mu_X \circ T\eta_X = \id_{TX}$, so it remains to check the other composite.
  By unitality and using $T\eta_X = \eta_{TX}$, we have
  \[
    \id_{T^2X} = \mu_{TX} \circ T\eta_{TX} = \mu_{TX} \circ T^2\eta_{X} .
  \]
  From the naturality square of $\mu$ applied to $\eta_X$, we get the following commutative square:
  \begin{center}
    \begin{tikzcd}
      T^2X \ar{d}{\mu_X} \ar{r}{T^2\eta_{X}} & T^3X \ar{d}{\mu_{TX}} \\
      TX \ar{r}{T\eta_{X}} & T^2X
    \end{tikzcd}
  \end{center}
  Combining these two facts, we deduce that $T\eta_X \circ \mu_X = \id_{T^2X}$, so $T\eta_X$ is an isomorphism.
\end{proof}


\begin{proposition}\label{prp:monad-factorization}
  Let $f : X \to Y$ be a morphism in $\mathcal{C}$. Then $f$ factors uniquely through $\eta_X$ if both $TX$ and $Y$ are $T$-reflexive.
\end{proposition}

\begin{proof}
  The naturality square of $\eta$ applied to $f$
\begin{center}
  \begin{tikzcd}
    X \ar{r}{f} \ar{d}{\eta_X} & Y \ar{d}{\eta_Y}[left]{\simeq} \\
    T X \ar{r}{T f} & T Y
  \end{tikzcd}
\end{center}
  yields a factorization $f = \eta_Y^{-1} \circ T f \circ \eta_X$ of $f$ through $\eta_X$.
  For uniqueness, assume we are given $f', f'' : TX \to Y$ such that $f' \circ \eta_X = f'' \circ \eta_X$.
  By \Cref{prp:monad-well-pointedness-torsionless-epi}, $\eta_X$ is $T$-torsionless-epic, which yields $f' = f''$.
\end{proof}

Now we specialize to a monad coming from an adjoint pair
$F \dashv G$ with $F : \mathcal{C} \to \mathcal{D}$, $G : \mathcal{D} \to \mathcal{C}$, unit $\eta : \id_{\mathcal{C}} \to GF$ and counit $\varepsilon : FG \to \id_{\mathcal{D}}$.
The associated monad is $T \coloneq GF$, with multiplication $\mu \coloneq G \varepsilon_F$ and unit $\eta$.
In the case where the monad is induced from an adjunction, we can make a slight improvement to \Cref{prp:monad-well-pointedness-torsionless-epi}.

\begin{proposition}\label{prp:monad-adjunction-well-pointed-improvement}
  Let $X$ be an object of $\mathcal{C}$.
  Then $GF\eta_X = \eta_{GFX}$ if and only if $F\eta_X$ is an isomorphism.
\end{proposition}

\begin{proof}
  Assume first that $GF\eta_X = \eta_{GFX}$.
  This part of the proof is completely analogous to the implication (ii) $\Rightarrow$ (i) in \Cref{prp:monad-well-pointedness-torsionless-epi}, but we write it out for completeness.
  From the triangle identity we know that $\varepsilon_{FX} \circ F\eta_X = \id_{FX}$.
  Consider the naturality square of $\varepsilon$ applied to $F\eta_X$:
  \begin{center}
    \begin{tikzcd}
      FGFX \ar{r}{FGF\eta_X} \ar{d}{\varepsilon_{FX}} & FGFGFX \ar{d}{\varepsilon_{FGFX}} \\
      FX \ar{r}{F\eta_X} & FGFX
    \end{tikzcd}
  \end{center}
  Combining this with the assumption $GF\eta_X = \eta_{GFX}$, we get that $F\eta_X \circ \varepsilon_{FX} = \varepsilon_{FGFX} \circ FGF\eta_X = \varepsilon_{FGFX} \circ F\eta_{GFX} = \id_{FGFX}$, where the last equality is given by the triangle identity.
  Thus $F\eta_X$ is an isomorphism, with inverse $\varepsilon_{FX}$.

  Now assume $F\eta_X$ is an isomorphism.
  By functoriality, $GF\eta_X$ is also an isomorphism with inverse $\mu_X$.
  Since $\eta_{GFX}$ also has $\mu_X$ as inverse, we get that $GF\eta_X = \eta_{GFX}$ by uniqueness of inverses.
\end{proof}

Let us apply the above discussion to the contravariant functor $(\argmhole)^\ast$.
This functor is right adjoint to itself, since $\Hom_R(A, B^\ast) \simeq \Hom_R(B,A^\ast)$.
To fit this into the above discussion, we take $F = (\argmhole)^\ast : \modcat{R} \to \modcat{R}^{\op}$
and $G = (\argmhole)^\ast : \modcat{R}^{\op} \to \modcat{R}$.
The unit and counit of this adjunction are both given by the evaluation map $\ev_A : A \to A^{\ast\ast}$.
Thus the associated monad is $T \coloneq (\argmhole)^{\ast\ast}$, with multiplication $\mu_A \coloneq (\ev_{A^\ast})^\ast$.
A module $X$ is $T$-reflexive if $\ev_X$ is an isomorphism, which is equivalent to $X$ being reflexive.
We say that a module $X$ is \notion{torsionless} if it is $T$-torsionless for this $T$,
i.e., if $\ev_X$ is a monomorphism.
This is equivalent to being torsionless in the sense of H. Bass~\cite{bass1960} and motivates the name.
We can specialize \Cref{prp:monad-well-pointedness-torsionless-epi} with the improvement from \Cref{prp:monad-adjunction-well-pointed-improvement} to this situation.

\begin{corollary}\label{cor:dbl-dual-reflexive-iss-dual-reflexive}
  Let $A$ be an $R$-module. The following are equivalent:
  \begin{enumerate}[noitemsep]
    \item $A^{\ast\ast}$ is reflexive
    \item $A^\ast$ is reflexive
    \item $\ev_A$ is torsionless-epic. \earlyqed
  \end{enumerate}
\end{corollary}

\begin{proposition}\label{prp:dbl-dual-tensor-factorization}
  Let $A$, $B$, and $C$ be finitely adequate modules. Then every map $f : A \otimes B \to C$ factors uniquely through $\ev_{A \otimes B} : A \otimes B \to A \dbldualtensor B$.
\end{proposition}

\begin{proof}
  By \Cref{prp:monad-factorization}, we see that $f$ factors uniquely through $\ev_{A \otimes B}$ if both $A \dbldualtensor B$ and $C$ are reflexive.
  The module $C$ is reflexive by \Cref{thm:can-dual-iso}, as it is finitely adequate.
  The double dual tensor product $A \dbldualtensor B$ is reflexive if $(A \otimes B)^\ast$ is reflexive by \Cref{cor:dbl-dual-reflexive-iss-dual-reflexive}.
  Now, $(A \otimes B)^\ast = \Hom_R(A \otimes B, R) \simeq \Hom_R(A, B^\ast)$, which is reflexive by \Cref{prp:fin-adeq-hom-closure} since $A$ and $B$ are finitely adequate.
\end{proof}

\begin{corollary}\label{cor:dbl-dual-tensor-adjunction}
  Let $M$ be a finitely adequate module.
  The functor $(\argmhole) \dbldualtensor M : \finadeq \to \finadeq$ is left adjoint to $\Hom_R(M,\argmhole) : \finadeq \to \finadeq$. \earlyqed
\end{corollary}


\printbibliography

\end{document}
